\scanpage{104}
% L1--6
\uline{Définition} 1 Soit $X$ un espace topologique. On appelle topologie constructible de $X$, \deleted{\unsure{et on}} la topologie dont les ouverts sont les réunions quelconques de parties loc. constructibles [donc une base d'ouverts \unsure{fournie|formée} par les parties loc constructibles,]. Parties \unsure{antitypiques|antypiques|indéfiques|analytiques} $=$ fermés de $X^{\mathrm{cons}}$ $=$ intersections des parties loc constructibles.

% L7--8; several separate vertical emphasis strokes at the left of the proposition.
\margin{\deleted{\illegible{2}}}
% slanted
\uline{Prop 1} Supposons que \deleted{tout point de} $X$ admette une base \inserted{$(U_i)_{i\in I}$} formée d'\deleted{ouverts} rétrocompacts. Alors

% L9
(0) les parties constructibles de $X$ forment une base d'ouverts de $X^{\mathrm{cons}}$

% L10--11
(i) La topologie constructible de $X$ est plus fine que la topologie donnée, \unsure{i.e.} $X^{\mathrm{cons}}\longrightarrow X$ est continue.

% L12--13
(ii) \deleted{\unsure{Si de|Supposons}} \inserted{Si} $X$ satisfait de plus $T_0$, alors $X^{\mathrm{cons}}$ est séparé [et $=$ tot. discontinu] et réciproquement

% L14--17
(iii) Si \deleted{$\unsure{U|E}$} $E$ est une \deleted{ouvert de $X$, la topologie} \inserted{partie constructible de $X$,} munie de la topologie induite $T_E$, alors la topologie constructible $T_E^{\mathrm{cons}}$ associée : $T_E$ est induite par $T^{\mathrm{cons}}$ sur $E$.

% L18--19
(0) Car \deleted{\unsure{tt}} un loc. constructible $E$ est réunion des $U_i\cap E$, qui sont constructibles dans $X$ par $\unsure{0_{\mathrm I}|0_{\mathrm{III}}}$ 9.1.8.

% L20--26
(i) Résulte du fait que les $U_i$ sont \inserted{[loc.]} constructibles, donc ouverts pour $T_{\unsure{E|X}}^{\mathrm{cons}}$. Pour (ii), \deleted{\unsure{soit|soient}}, $x,y\in X$, il existe \unsure{un des deux|l'un des deux}, disons $x$, et un voisinage ouvert $U\ni x$, $U\not\ni y$. On peut grâce à l'hyp supposer $U$ rétrocomp\unsure{act}, donc [loc.] constructible, et $X-U=V$ l'est \inserted{alors} également, alors $U\cap V=\emptyset$, $U\cup V=X$, $U\ni x$, $V\ni y$, $U,V$ ouverts pour $T_E$.

% L27--30
Pour (iii), on note que la topologie induite par $T_{\unsure{X|E}}^{\mathrm{cons}}$ sur $E$ a comme base d'ouverts les $E\cap F$, $F$ \deleted{\unsure{un|const.}} \inserted{loc.} constructible dans $X$, ce sont exactement les $G\subset E$ loc. constructibles dans $E$,
% archive sheet number: 103
